\documentclass{article}

\usepackage{fancyhdr}
\usepackage{amsmath}
\usepackage{amsthm}
\usepackage{amsfonts}
\usepackage{tikz}
\usepackage[plain]{algorithm}
\usepackage{algpseudocode}
\usepackage{amssymb}
\usepackage{enumitem}

\usetikzlibrary{automata,positioning}

\topmargin=-0.45in
\evensidemargin=0in
\oddsidemargin=0in
\textwidth=6.5in
\textheight=9.0in
\headsep=0.25in

\linespread{1.1}

\pagestyle{fancy}
\lhead{\hmwkAuthorName}
\chead{}
\rhead{\hmwkHeaderTitle}
\cfoot{\thepage}

\renewcommand\headrulewidth{0.4pt}
\renewcommand\footrulewidth{0pt}

\setlength\parindent{0pt}

\newcommand{\hmwkTitle}{Homework\ 6}
\newcommand{\hmwkDueDate}{Thursday, March 12, 2026,\ 11{:}59 pm (Thai time)}
\newcommand{\hmwkClass}{Data Structure\ and Abstractions}
\newcommand{\hmwkClassTime}{Section 2}
\newcommand{\hmwkClassInstructor}{Kanat Tangwongsan}
\newcommand{\hmwkAuthorName}{\textbf{Jiraroj Wiruchpongsanon (6781617)}}

\newcommand{\hmwkHeaderTitle}{Data Structure and Abstractions: Homework 6}

\title{
    \vspace{2in}
    \textmd{\textbf{\hmwkClass:\ \hmwkTitle}}\\
    \normalsize\vspace{0.1in}\small{Due\ on\ \hmwkDueDate}\\
    \vspace{0.1in}\large{\textit{\hmwkClassInstructor\ \hmwkClassTime}}
    \vspace{3in}
}

\author{\hmwkAuthorName}
\date{}

\renewcommand{\part}[1]{\textbf{\large Part \Alph{partCounter}}\stepcounter{partCounter}\\}

\newcounter{partCounter}
\newcommand{\solution}{\textbf{\large Solution}}

\newtheorem*{theorem}{Theorem}

\theoremstyle{remark}
\newtheorem*{remark}{Remark}

\begin{document}

\maketitle
\begin{center}
    \textit{Data Struct. and \{Abstractions, OOP\} (Term II/2025--26)}\\
    \textit{built on 2026/03/03 at 12:44:44}\\
    \textit{due: thu mar 12 @ 11:59pm}
\end{center}
\newpage

\section*{Task 4: Quick Sort Recurrence (4 points)}

Consider the recurrence
\[
    f(n) = n + 1 + \frac{2}{n} \bigl(f(n - 1) + f(n - 2) + \cdots + f(1)\bigr),
\]
where $f(0) = 0$.

This recurrence represents the ``average'' running time of the randomized quick sort algorithm\footnote{Technically, for those who took Discrete Math already, this is the expected running time.}. For now, we won't discuss how one can come up with such a recurrence. In this exercise, we're interested in solving this recurrence for a closed form. You'll do this in a few steps:

\begin{enumerate}
    \item[(i)] Consider $f(n)$ and $f(n - 1)$, where $n \ge 2$. We have
    \[
        f(n) = (n + 1) + \frac{2}{n} \bigl(f(n - 1) + f(n - 2) + \cdots + f(1)\bigr)
    \]
    and
    \[
        f(n - 1) = n + \frac{2}{n - 1} \bigl(f(n - 2) + f(n - 3) + \cdots + f(1)\bigr).
    \]
    By multiplying the first equation by $n$ and the second equation by $(n - 1)$, we have
    \begin{align}
        n \cdot f(n) &= (n + 1)n + 2 f(n - 1) + f(n - 2) + \cdots + f(1) \tag{1} \\
        (n - 1) f(n - 1) &= n(n - 1) + 2 f(n - 2) + f(n - 3) + \cdots + f(1). \tag{2}
    \end{align}
    Subtracting equation (2) from equation (1), we'll get
    \[
        n \cdot f(n) - (n - 1) f(n - 1) = 2n + 2 f(n - 1).
    \]
    In other words,
    \[
        n \cdot f(n) = 2n + (n + 1) f(n - 1). \tag{3}
    \]
    Your task in this step is to understand the derivation we just made. Other than that, no actions are required on your part.

    \item[(ii)] Let $g(n) = \frac{f(n)}{n + 1}$. Can you write what you have in terms of the function $g$? (Hint: divide equation (3) by $n(n + 1)$.)

    \item[(iii)] Your task in this step is to find a closed form for $g$. The recurrence $g$ that you have isn't listed in the table. But you can easily solve for a closed form. Before you begin, let us show you how to solve a related recurrence: $h(n) = h(n - 1) + \frac{1}{n}$, with $h(0) = 0$.

    We start out by unraveling $h(n)$. By the definition of $h(n)$,
    \[
        h(n) = h(n - 1) + \frac{1}{n}
    \]
    \[
        = h(n - 2) + \frac{1}{n} + \frac{1}{n - 1}
    \]
    \[
        = h(n - 3) + \frac{1}{n} + \frac{1}{n - 1} + \frac{1}{n - 2}
    \]
    expand the recurrence one more time

    expand the recurrence one more time

    It is pretty clear that if we keep on expanding the recurrence, we'll get
    \[
        h(n) = \frac{1}{n} + \frac{1}{n - 1} + \frac{1}{n - 2} + \cdots + \frac{1}{3} + \frac{1}{2} + 1.
    \]

    In Math, this has a name: the $n$-th Harmonic number, denoted by $H_n$, is given by $H_n = \frac{1}{n} + \frac{1}{n - 1} + \frac{1}{n - 2} + \cdots + \frac{1}{3} + \frac{1}{2} + 1$, so in terms of Harmonic number $h(n) = H_n$.

    \item[(iv)] Now that you can express $g(n)$ in terms of some expression involving Harmonic numbers, you can proceed to derive a closed form for $f(n)$. Finally, use the following fact to conclude that $f(n)$ is $O(n \ln n)$:

    Fact: $H_n \le 1 + \ln(n)$, where $\ln$ denotes the natural logarithm.
\end{enumerate}

\medskip
\solution

\begin{enumerate}[label=(\roman*)]
    \item at the end of the reading, we got
        \begin{equation}
            n \cdot f(n) - (n-1)f(n-1) = 2n + 2f(n-1)
        \tag{3}
        \end{equation}

    \item let

        \[g(n) = \frac{f(n)}{n+1}\]

        as hinted, we do 

        \begin{align*}
            \frac{(3)}{n(n+1)}: \frac{n \cdot f(n)}{n(n+1)} &= \frac{2n+(n+1)f(n-1)}{n(n+1)} \\
            g(n) &= \frac{f(n)}{n+1} \\
            g(n) &= \frac{2}{n+1} + \frac{f(n-1)}{n} && \quad \text{ as } g(n-1) = \frac{f(n-1)}{n}\\
        \end{align*}

        therefore, 

        \[            
            \boxed{g(n) = \frac{2}{n+1} + g(n-1)}
        \]

    \item write $g(n)$ out as

        \begin{align*}
            g(n) &= g(n-1) + \frac{2}{n+1} \\
            &= g(n-2) + \frac{2}{n} + \frac{2}{n+1} \\
            &= g(n-3) + \frac{2}{n-1} + \frac{2}{n} + \frac{2}{n+1} \\
            &... \\
            g(n) &= \frac{2}{0+1} + \frac{2}{1+1} + \frac{2}{2+1} + ... + \frac{2}{n-1} + \frac{2}{n} + \frac{2}{n+1}
        \end{align*}

        so, 

        \[\boxed{g(n) = \sum_{i=0}^n \frac{2}{n+1}}\]

        and as we're given a \textbf{fact} that

        \[H_n \le 1 + \ln{n}\]

        we got

        \[\boxed{g(n) = 2H_{n+1}}\]
\end{enumerate}

\end{document}
